
\section{Introduction}

Phishing combines social engineering and technical deception
to impersonate legitimate services and steal sensitive
information through websites, email, and SMS. Modern attacks
coordinate deceptive information across URLs, HTML, rendered
pages, messages, and hosting infrastructure, making
single-source detection insufficient. The challenge is
particularly severe for previously unseen attacks, which
cannot be identified through existing reputation records.
Lin et al.~\cite{lin2021phishpedia} identified 1,133 phishing
pages undetected by any VirusTotal engine among 1,704
collected pages, with 74.6\% remaining unreported after
one week. Chiba et al.~\cite{chiba2026phishlumos} reported
discovering previously unlisted phishing URLs a median
of 192.8 hours before expert verification.

Consider a phishing SMS directing recipients to a
credential-harvesting page hosted on a legitimate cloud
platform. The message and page display the same
impersonated brand, while the URL appears ordinary and
the hosting platform has a valid TLS certificate and a
long registration history. Reputation-based indicators
may therefore be unavailable or misleading. Even a
multimodal detector may incorrectly interpret the
repeated brand claim as independent corroboration,
although both observations originate from a single
attacker-controlled choice. Conversely, conflicting
observations across modalities may reveal deception
that individual agents cannot identify independently.
Reliable detection thus requires item-level provenance,
common-cause-aware evidence fusion, and explicit
resolution of conflicting findings.

Existing machine learning (ML)-based phishing detectors
rely heavily on predefined features and labeled training
data, limiting their ability to generalize to previously
unseen attacks and evolving phishing strategies.
Although large language models (LLMs) offer greater
flexibility in interpreting heterogeneous evidence,
their reasoning may be inconsistent, and self-reported
confidence is often unreliable~\cite{xiong2024can}.
These limitations are further exacerbated by
distribution shift, which degrades calibration and
the reliability of confidence-based
decisions~\cite{ovadia2019can}. Recent multi-agent
approaches employ modality-specialized analysis,
but agreement among agents may reflect correlated
evidence rather than independent corroboration.
Moreover, confidence-based escalation can be
ineffective when specialists have heterogeneous
capabilities~\cite{jitkrittum2023when}, while
repeatedly invoking all agents introduces
unnecessary computational overhead.

Collaborative reasoning introduces further challenges,
as agents may adopt peer conclusions without sufficient
supporting evidence, making their decisions difficult
to verify. Furthermore, conventional evaluations may
overestimate zero-day detection capabilities because
test samples absent from experimental training data
may still have appeared in an LLM's pretraining corpus.
These limitations highlight the need for
provenance-aware evidence fusion, reliable uncertainty
estimation, selective and verifiable collaborative
reasoning, and rigorous zero-day evaluation that
accounts for both reputation-source absence and
model knowledge cutoffs.



To address these challenges, we propose
\textbf{MA-ZeroPhish}, an adaptive agentic
multi-agent framework for zero-day phishing
detection. It integrates shared evidence
acquisition, five modality-specialized agents,
provenance-aware collaboration, and independent
adjudication. An Orchestrator selects specialists
using structural evidence conditions, while a
Moderator reconciles dependent observations
and initiates targeted investigation based on
estimated stopping error and actionable issues.
An independent Judge produces evidence-supported
verdicts without access to specialist judgments
or confidence bands. The main contributions
are as follows:

\begin{enumerate}

\item \textbf{Adaptive evidence acquisition
and specialist orchestration.}
We develop a four-phase architecture integrating
shared acquisition with five specialists for
URL, web-structure, rendered-content, SMS/email,
and infrastructure-metadata analysis.
Structural triggers guide budget-constrained
specialist selection, while declared tool
predicates govern supplementary acquisition,
enabling auditable, case-specific investigation.

\item \textbf{Provenance-aware reasoning
and uncertainty-driven collaboration.}
We introduce validated specialist findings,
deterministic evidence-based confidence bands,
and common-cause-aware reconciliation to prevent
correlated observations from being counted as
independent support. A calibrated stopping-error
estimator triggers targeted collaboration only
when actionable evidence issues remain, subject
to budget and round limits.

\item \textbf{Auditable collaboration
and independent adjudication.}
We design evidence-linked collaboration with
attributed conflict resolution, blinded
investigation, and verifiable predecessor-linked
revisions. An independent Judge evaluates
eligible observations, provenance dependencies,
and coverage limitations without specialist
verdicts or confidence bands. We further
specify a zero-day evaluation methodology
using temporal and configuration-data
separation, reputation-source exclusion,
and targeted ablations to assess detection
performance and investigation cost.

\end{enumerate}




